\begin{figure*}[t]
    \centering
    \includegraphics[width=\textwidth]{figures/main/methods_pipeline/methods_pipeline.pdf}
    \caption{Overview of \methodname{}. Calibration-guided geometry constructs
    canonical frame transforms from image-plane metadata and a regular-sweep
    prior. A shared 3D Gaussian scene is rendered by complementary local
    finite-slice and global transmission pathways. The resulting image
    residuals and renderer-specific contribution statistics guide scene
    optimization and Gaussian adaptation.}
    \label{fig:methods-pipeline}
\end{figure*}

\section{Method}
\label{sec:method}

In this section, we first formulate Gaussian-based sensorless 3D ultrasound
reconstruction and define the Gaussian scene representation
(Sec.~\ref{sec:problem-formulation-and-scene-representation}). Subsequently,
we introduce an acquisition-guided pseudo-pose parameterization to recover
coherent scan geometry without measured probe poses
(Sec.~\ref{sec:acquisition-guided-pseudo-pose-parameterization}). 
Next, we develop a physics-aware
local-global Gaussian renderer for physically faithful B-mode synthesis
(Sec.~\ref{sec:physics-aware-local-global-gaussian-renderer}). 
Furthermore, contribution-aware
Gaussian adaptation is proposed to regulate the Gaussian population for
ultrasound reconstruction
(Sec.~\ref{sec:contribution-aware-gaussian-adaptation}). 
The overall framework and key module designs are illustrated in
Fig.~\ref{fig:methods-pipeline}.

\vspace{-3mm}
\subsection{Problem Formulation and Scene Representation}
\label{sec:problem-formulation-and-scene-representation}

We consider Gaussian-based sensorless 3D ultrasound reconstruction from an
ordered sequence of $N$ B-mode frames,
$\mathcal I=\{\mathbf I_i\}_{i=1}^{N}$, where the corresponding measured
probe-pose sequence
$\mathcal T=\{\mathbf T_i^{p\rightarrow w}\}_{i=1}^{N}$ is unavailable.
Instead, we associate the frames with a computationally derived pseudo-pose
sequence
$\widetilde{\mathcal T}
=
\{\widetilde{\mathbf T}_i^{p\rightarrow w}\}_{i=1}^{N}$,
where each transform maps the probe coordinate system of frame $i$ to a shared
world frame. Given the image sequence and ultrasound
calibration, our goal is to recover a consistent 3D ultrasound representation
whose rendered B-mode observations agree with the frames.

Following 3D Gaussian Splatting (3DGS)
\cite{kerbl3DGaussianSplatting2023}, we represent the ultrasound scene as a
set of $G$ anisotropic Gaussian primitives,
$\mathcal G=\{\mathbf G_g\}_{g=1}^{G}$, while preserving the same Gaussian
geometry. Each Gaussian is parameterized by its center
$\boldsymbol\mu_g\in\mathbb R^3$, anisotropic scale
$\mathbf s_g\in\mathbb R_+^3$, unit quaternion $\mathbf q_g$, echo attribute
$\alpha_g$, and transmission attribute $\tau_g$. Its covariance is derived as
$\boldsymbol\Sigma_g
=
\mathbf R_g\mathbf S_g\mathbf S_g^\top\mathbf R_g^\top$,
where $\mathbf R_g=\mathbf R(\mathbf q_g)$ and
$\mathbf S_g=\operatorname{diag}(\mathbf s_g)$. Unlike conventional 3DGS,
which uses opacity and spherical-harmonic color coefficients to model optical
appearance, $\alpha_g$ and $\tau_g$ characterize local echo formation and
global acoustic transmission, respectively, 
thereby enabling more physically faithful B-mode image synthesis.

Given the ordered B-mode sequence $\mathcal I$, pseudo-pose sequence
$\widetilde{\mathcal T}$, and an initialized Gaussian scene $\mathcal G$, we
render each frame as
$\widehat{\mathbf I}_i
=
\mathcal R(\mathcal G,\widetilde{\mathbf T}_i^{p\rightarrow w})$,
where $\mathcal R$ denotes the ultrasound-specific Gaussian renderer. We then
minimize the reconstruction loss $\mathcal L$ between
$\widehat{\mathbf I}_i$ and its corresponding observation $\mathbf I_i$ to
iteratively optimize $\mathcal G$, while periodically updating the Gaussian
population through the adaptation operator
$\mathcal G\leftarrow\mathcal A(\mathcal G)$, ultimately yielding the
optimized scene $\mathcal G^{*}$:
\begin{equation}
\mathcal G^{*}
=
\arg\min_{\mathcal G}
\sum_{i=1}^{N}
\mathcal L
\left(
\widehat{\mathbf I}_i,
\mathbf I_i
\right).
\label{eq:overall-reconstruction}
\end{equation}

Similar to 3DGS, we instantiate $\mathcal L$ as a weighted combination of
$\ell_1$ and SSIM terms:
\begin{equation}
\mathcal L
=
(1-\lambda)
\left\|
\widehat{\mathbf I}_i-\mathbf I_i
\right\|_1
+
\lambda
\left[
1-
\operatorname{SSIM}
\left(
\widehat{\mathbf I}_i,
\mathbf I_i
\right)
\right].
\label{eq:total-loss}
\end{equation}
Here, $\lambda=0.2$ balances image fidelity and structural consistency during
optimization.

\vspace{-3mm}
\subsection{Acquisition-Guided Pseudo-Pose Parameterization}
\label{sec:acquisition-guided-pseudo-pose-parameterization}

Existing sensorless ultrasound reconstruction methods typically recover
frame-wise probe motion from adjacent B-mode images using learned pose
estimators, which not only require additional training and computation but
also suffer from accumulated errors when pairwise estimates are chained.
Moreover, reliable motion inference from ultrasound appearance remains
challenging due to limited inter-slice overlap, repetitive anatomy,
view-dependent speckle, and substantial intensity variation. In contrast, we
observe that practical ultrasound sweeps often exhibit strong acquisition
regularity, with approximately parallel imaging planes and
translation-dominant probe motion rather than large frame-to-frame rotations.
Motivated by this observation, we replace frame-wise 6-DoF motion estimation
with an acquisition-guided pseudo-pose parameterization that exploits frame
order and nominal scan progression to efficiently construct the pseudo-pose
sequence $\widetilde{\mathcal T}$ used by ULGSplat.

Specifically, given the translation-dominant acquisition pattern, we model the
inter-frame progression by a shared effective spacing $d=\rho d_0$, where
$d_0>0$ is the nominal acquisition spacing and $\rho>0$ globally adjusts 
the effective scan spacing. Without loss of generality, we take the probe coordinate system of
a reference frame $i_0$ as the canonical world frame and use a predefined 
signed canonical axis $\mathbf u$ that specifies the scan direction.
Accordingly, rather than estimating frame-wise motion independently, the
probe-to-world pseudo pose of frame $i$ is directly induced by its relative
frame index:
\begin{equation}
\widetilde{\mathbf T}_i^{p\rightarrow w}
=
\begin{bmatrix}
\mathbf I_3 & (i-i_0)d\mathbf u\\
\mathbf 0^\top & 1
\end{bmatrix}.
\label{eq:calibration-geometry-single}
\end{equation}

This shared acquisition model yields a geometrically consistent pseudo-pose
sequence in three respects.
First, mapping every frame to the same canonical world frame provides a unified
spatial reference. A point $\bar{\mathbf x}^{p_i}$ in frame $i$ is mapped as
$\bar{\mathbf x}^{w}
=
\widetilde{\mathbf T}_i^{p\rightarrow w}\bar{\mathbf x}^{p_i}$,
allowing observations from different B-mode slices to be organized directly
within the shared Gaussian scene.
Furthermore, the relative transform between frames $i$ and $j$ is analytically
determined by the shared trajectory:
\begin{equation}
\widetilde{\mathbf T}_{j\leftarrow i}
=
\left(
\widetilde{\mathbf T}_j^{p\rightarrow w}
\right)^{-1}
\widetilde{\mathbf T}_i^{p\rightarrow w}
=
\begin{bmatrix}
\mathbf I_3 & (i-j)d\mathbf u\\
\mathbf 0^\top & 1
\end{bmatrix},
\label{eq:pseudo-pose-relative}
\end{equation}
which naturally guarantees bidirectional consistency,
$\widetilde{\mathbf T}_{i\leftarrow j}
=
\widetilde{\mathbf T}_{j\leftarrow i}^{-1}$.
More importantly, the transforms satisfy composition consistency
$\widetilde{\mathbf T}_{k\leftarrow j}
\widetilde{\mathbf T}_{j\leftarrow i}
=
\widetilde{\mathbf T}_{k\leftarrow i}$
for arbitrary frames $i$, $j$, and $k$, ensuring that their relative geometry
is independent of the transformation path. 

Consequently, by combining acquisition regularity with shared trajectory
parameterization, the proposed method simplifies pseudo-pose construction
while preserving geometric consistency, eliminating the need for
appearance-based frame-wise motion inference and pairwise motion chaining, and
providing a reliable spatial basis for subsequent Gaussian initialization and
rendering.

\vspace{-2mm}
\subsection{Physics-Aware Local-Global Gaussian Renderer}
\label{sec:physics-aware-local-global-gaussian-renderer}

Conventional 3DGS renders an image by projecting 3D Gaussian primitives onto
the image plane and compositing their opacity-weighted appearance in
front-to-back depth order. Nevertheless, B-mode ultrasound follows acoustic
focusing within a finite elevational slice, so each pixel may aggregate echoes
from multiple scatterers within the local acoustic support; meanwhile, echoes
from deeper structures are modulated by tissues encountered earlier along the
propagation path, making the observed intensity jointly dependent on
finite-slice echo formation and preceding-path acoustic propagation.
To address this mismatch, we develop a physics-aware local-global Gaussian
renderer in ULGSplat, comprising a finite-slice local echo pathway and a
path-dependent global transmission pathway. The local pathway models the
finite-slice local echo response by accounting for depth-dependent slice
sensitivity and the in-plane geometry of scatterers within the acoustic slice,
whereas the global pathway models the path-dependent transmission response by
tracing the acoustic propagation path and accumulating the effects of causally
preceding structures. These two responses are subsequently composed to
synthesize physically faithful B-mode images.

\subsubsection{Finite-Slice Local Echo Formation}

To preserve spatial correspondence between pixels and scatterers
within the finite acoustic slice, we redesign the Gaussian attribute derivation
of the local pathway with two key components: depth-adaptive elevational
conditioning, which models depth-dependent slice sensitivity, and probe-plane
Gaussian marginalization, which maintains image-plane Gaussian geometry.

We first establish the probe-space representation shared by these operations.
Given the pseudo pose
$\widetilde{\mathbf T}_i^{p\rightarrow w}
=[\mathbf R_i\mid\mathbf t_i]$,
Gaussian $g$ is transformed into probe space as
\begin{equation}
\boldsymbol\mu^{p}_{ig}
=
\mathbf R_i^\top
\left(
\boldsymbol\mu_g-\mathbf t_i
\right),
\qquad
\boldsymbol\Sigma^{p}_{ig}
=
\mathbf R_i^\top
\boldsymbol\Sigma_g
\mathbf R_i,
\label{eq:probe-space-gaussian}
\end{equation}
yielding its probe-space mean
$\boldsymbol\mu^{p}_{ig}$ and covariance
$\boldsymbol\Sigma^{p}_{ig}$, which are also reused by the global transmission
pathway.

On this probe-space representation, we introduce a depth-adaptive elevational
kernel centered on the imaging plane to model the finite-slice geometry and
echo formation of each primitive. Given the calibrated axial direction
$\mathbf e_{\mathrm{ax}}$ and near depth $z_{\mathrm{near}}$, the axial depth
of Gaussian $g$ and the corresponding elevational kernel width are given by
\begin{equation}
\begin{aligned}
z_{ig}
&=
\max
\left(
\mathbf e_{\mathrm{ax}}^\top
\boldsymbol\mu^p_{ig}
-
z_{\mathrm{near}},
0
\right),\\
\sigma^{\mathrm{ele}}_{ig}
&=
\operatorname{clip}
\left(
\sigma^{\mathrm{ele}}_0
+
\sigma^{\mathrm{ele}}_1 z_{ig},
\sigma^{\mathrm{ele}}_{\min},
\sigma^{\mathrm{ele}}_{\max}
\right),
\end{aligned}
\label{eq:depth-adaptive-ele-width}
\end{equation}
where $z_{ig}$ denotes the axial depth relative to the calibrated near plane,
$\sigma^{\mathrm{ele}}_0$ and $\sigma^{\mathrm{ele}}_1$ control the initial
width and its depth-dependent growth, and
$\sigma^{\mathrm{ele}}_{\min}$ and $\sigma^{\mathrm{ele}}_{\max}$ bound the
effective elevational support. Specifically, multiplying this kernel by the
probe-space Gaussian
$(\boldsymbol\mu^{p}_{ig},\boldsymbol\Sigma^{p}_{ig})$, with
$\boldsymbol\Lambda^{p}_{ig}
=(\boldsymbol\Sigma^{p}_{ig})^{-1}$, yields

\begin{equation}
\begin{aligned}
\boldsymbol\Sigma^{e}_{ig}
&=
\left[
\boldsymbol\Lambda^{p}_{ig}
+
\left(\sigma^{\mathrm{ele}}_{ig}\right)^{-2}
\mathbf e_{\mathrm{ele}}
\mathbf e_{\mathrm{ele}}^\top
\right]^{-1},
\qquad
\boldsymbol\mu^{e}_{ig}
=
\boldsymbol\Sigma^{e}_{ig}
\boldsymbol\Lambda^{p}_{ig}
\boldsymbol\mu^{p}_{ig},
\\
r^{\mathrm{ele}}_{ig}
&=
\frac{
\left(
\mathbf e_{\mathrm{ele}}^\top
\boldsymbol\mu^{p}_{ig}
\right)^2
}{
\mathbf e_{\mathrm{ele}}^\top
\boldsymbol\Sigma^{p}_{ig}
\mathbf e_{\mathrm{ele}}
+
\left(\sigma^{\mathrm{ele}}_{ig}\right)^2
+
\varepsilon
},
\qquad
\eta_{ig}
=
\exp
\left(
-\frac{1}{2}r^{\mathrm{ele}}_{ig}
\right),
\end{aligned}
\label{eq:elevational-conditioned-attributes}
\end{equation}
where $\boldsymbol\mu^{e}_{ig}$ and
$\boldsymbol\Sigma^{e}_{ig}$ denote the conditioned mean and covariance,
while $r^{\mathrm{ele}}_{ig}$ and $\eta_{ig}$ denote the normalized
elevational distance and the resulting slice-overlap factor, respectively.

With the conditioned mean $\boldsymbol\mu^{e}_{ig}$ and covariance
$\boldsymbol\Sigma^{e}_{ig}$, we apply probe-plane Gaussian
marginalization to derive the image-plane Gaussian attributes for local echo
formation, together with a footprint-area factor for response compensation.
Given the calibrated lateral and axial directions
$\mathbf e_{\mathrm{lat}}$ and $\mathbf e_{\mathrm{ax}}$, and the associated
projection matrix
$\mathbf P_{\mathrm{la}}
=
[\mathbf e_{\mathrm{lat}},
\mathbf e_{\mathrm{ax}}]^\top$,
we obtain
\begin{equation}
\begin{aligned}
\boldsymbol\mu^{\mathrm{loc},2D}_{ig}
&=
\mathbf P_{\mathrm{la}}
\boldsymbol\mu^{e}_{ig},
\qquad
\boldsymbol\Sigma^{\mathrm{loc},2D}_{ig}
=
\mathbf P_{\mathrm{la}}
\boldsymbol\Sigma^{e}_{ig}
\mathbf P_{\mathrm{la}}^\top,
\\
\nu_{ig}
&=
\left[
\det
\left(
\boldsymbol\Sigma^{\mathrm{loc},2D}_{ig}
\right)
+
\varepsilon
\right]^{-\gamma},
\end{aligned}
\label{eq:local-image-attributes}
\end{equation}
where $\boldsymbol\mu^{\mathrm{loc},2D}_{ig}$ and
$\boldsymbol\Sigma^{\mathrm{loc},2D}_{ig}$ denote the image-plane mean and
covariance for local echo formation, $\nu_{ig}$ is the footprint-area factor,
and $\gamma$ controls the degree of area compensation.

To enable efficient rendering, we construct a sparse local pixel--Gaussian
candidate set based on the image-space geometric support of each Gaussian and
its influence on the queried pixel. Similar to vanilla 3DGS, 
we formulate the image-space influence of Gaussian
$g$ on pixel $p$ as
\begin{equation}
\phi_{igp}
=
\exp
\left[
-\frac{1}{2}
\mathbf d_{igp}^{\top}
\left(
\boldsymbol\Sigma^{\mathrm{loc},2D}_{ig}
\right)^{-1}
\mathbf d_{igp}
\right],
\label{eq:local-gaussian-response}
\end{equation}
where
$\mathbf d_{igp}
=
\boldsymbol\xi^{2D}_{p}
-
\boldsymbol\mu^{\mathrm{loc},2D}_{ig}$
denotes the image-plane offset from Gaussian $g$ to pixel $p$.
The local candidate set is then defined as
\begin{equation}
\mathcal N^{\mathrm{loc}}_{ip}
=
\left\{
g
\;\middle|\;
\boldsymbol\xi^{2D}_{p}
\in
\mathcal B^{\mathrm{loc}}_{ig},
\;
\eta_{ig}\nu_{ig}\phi_{igp}
>
\tau_{\mathrm{loc}}
\right\},
\label{eq:local-candidate-set}
\end{equation}
where $\mathcal B^{\mathrm{loc}}_{ig}$ denotes the compact image-space support
of Gaussian $g$, and $\tau_{\mathrm{loc}}$ is a small threshold for discarding
negligible contributions.

Finally, the local echo response is obtained by additive accumulation over the
sparse candidate set:
\begin{equation}
L_{ip}
=
\sum_{g\in\mathcal N^{\mathrm{loc}}_{ip}}
\alpha_g
\eta_{ig}
\nu_{ig}
\phi_{igp},
\label{eq:local-response}
\end{equation}
where $\alpha_g$ is the learnable local echo amplitude and $L_{ip}$ is the
accumulated local echo response at pixel $p$ in frame $i$.

\subsubsection{Path-Dependent Global Transmission}

Using the local pathway alone is insufficient to capture ultrasound image
formation, as it ignores structures along the acoustic path and
therefore cannot model path-dependent attenuation and transmission. To address
this limitation, we introduce path-dependent global transmission through two
components: ray-aware pixel--Gaussian association and causal ray integration.

To identify the Gaussians that can causally affect pixel $p$, we construct a
sparse global pixel--Gaussian candidate set from the interaction between each
probe-space Gaussian and the queried pixel's propagation path. Given the
calibrated pixel coordinate
$\boldsymbol\xi^{2D}_{p}
=
[\xi_{\mathrm{lat},p},\xi_{\mathrm{ax},p}]^\top$,
the corresponding pre-pixel acoustic ray in probe space is given by
\begin{equation}
\mathbf q_{ip}(t)
=
\xi_{\mathrm{lat},p}\mathbf e_{\mathrm{lat}}
+
\left(z_{\mathrm{near}}+t\right)\mathbf e_{\mathrm{ax}},
\qquad
0\leq t\leq t_p,
\label{eq:global-ray}
\end{equation}
where
$t_p=\xi_{\mathrm{ax},p}-z_{\mathrm{near}}$
denotes the propagation distance from the calibrated near plane to the target
pixel. Based on this ray, we retain only Gaussians that intersect the
propagation segment and satisfy depth causality, yielding the candidate set:
\begin{equation}
\mathcal N^{\mathrm{glo}}_{ip}
=
\left\{
g
\;\middle|\;
\mathcal B^{\mathrm{glo}}_{ig}
\cap
\mathcal R_{ip}
\neq
\varnothing,\;
\mathbf e_{\mathrm{ax}}^\top
\boldsymbol\mu^{p}_{ig}
\leq
\xi_{\mathrm{ax},p}
\right\},
\label{eq:global-candidate-set}
\end{equation}
where
$\mathcal R_{ip}
=
\{\mathbf q_{ip}(t)\mid0\leq t\leq t_p\}$
denotes the pre-pixel acoustic segment, and
$\mathcal B^{\mathrm{glo}}_{ig}$ denotes the compact probe-space support of
Gaussian $g$ used for ray association.

For each candidate Gaussian, we evaluate its influence along the acoustic path
of pixel $p$. Specifically, the pointwise influence of Gaussian $g$ at
$\mathbf q_{ip}(t)$ is formulated as
\begin{equation}
\psi_{igp}(t)
=
\exp
\left[
-\frac{1}{2}
\mathbf d_{igp}(t)^\top
\boldsymbol\Lambda^{p}_{ig}
\mathbf d_{igp}(t)
\right],
\label{eq:global-ray-influence}
\end{equation}
where
$\mathbf d_{igp}(t)
=
\mathbf q_{ip}(t)-\boldsymbol\mu^{p}_{ig}$
denotes the probe-space offset from the Gaussian center to the queried point
on the acoustic ray. The ideal ray--Gaussian overlap is defined by
\begin{equation}
\bar{\rho}_{igp}
=
\int_{0}^{t_p}
\psi_{igp}(t)\,dt,
\label{eq:causal-ray-overlap-ideal}
\end{equation}
which accumulates the pointwise influence over the complete pre-pixel
propagation path.

Directly evaluating Eq.~\eqref{eq:causal-ray-overlap-ideal} over the full
pre-pixel interval for every ray--Gaussian pair is computationally expensive.
Since the Gaussian influence rapidly decays away from its closest interaction
with the ray, we restrict numerical integration to a compact neighborhood of
this location. Specifically, we first locate the point of strongest
ray--Gaussian interaction and then expand a compact integration interval around
it:
\begin{equation}
\begin{aligned}
t^{*}_{igp}
&=
\operatorname*{arg\,min}_{0\leq t\leq t_p}
\mathbf d_{igp}(t)^\top
\boldsymbol\Lambda^{p}_{ig}
\mathbf d_{igp}(t),\\
t_s
&=
\max(0,t^{*}_{igp}-w),
\qquad
t_e
=
\min(t_p,t^{*}_{igp}+w),
\end{aligned}
\label{eq:localized-ray-interval}
\end{equation}
where $t^{*}_{igp}$ is the interaction center, $w$ the integration half-window,
and $t_s,t_e$ the clipped endpoints. Integrating over this interval yields 
the approximated ray--Gaussian overlap using low-order Gauss--Legendre quadrature:
\begin{equation}
\rho_{igp}
=
\int_{t_s}^{t_e}
\psi_{igp}(t)\,dt
\approx
\bar{\rho}_{igp}.
\label{eq:localized-ray-overlap}
\end{equation}

Using $\rho_{igp}$, we convert each candidate Gaussian into a transmission
factor through its learnable transmission floor $\tau_g$, and multiplicatively
accumulate these factors:
\begin{equation}
\begin{aligned}
T_{igp}
&=
\tau_g
+
\left(
1-\tau_g
\right)
\exp
\left(
-\lambda_{\mathrm{ray}}\rho_{igp}
\right),\\
T_{ip}
&=
\operatorname{clip}
\left(
\prod_{g\in\mathcal N^{\mathrm{glo}}_{ip}}
T_{igp},
T_{\min},
1
\right),
\end{aligned}
\label{eq:global-transmission}
\end{equation}
where $\lambda_{\mathrm{ray}}$ controls the attenuation strength,
$T_{igp}$ is the per-Gaussian transmission factor, and $T_{\min}$ prevents
collapse of the accumulated global transmission response $T_{ip}$.

\subsubsection{Local-Global Response Composition}

Finally, we use the global transmission response $T_{ip}$ to modulate the local
echo response $L_{ip}$, where $T_{ip}$ is first blended with identity
transmission through $\beta$ to regulate how strongly path-dependent global
transmission affects finite-slice local echo formation:
\begin{equation}
\widehat I_{ip}
=
L_{ip}
\left(
1-\beta+\beta T_{ip}
\right),
\label{eq:local-global-composition}
\end{equation}
where $\beta\in[0,1]$ controls the strength of global transmission modulation,
with $\beta=0$ retaining the local echo response and $\beta=1$ applying the
full transmission response. 

By coupling spatially localized echo formation with preceding-path acoustic
modulation, the proposed renderer preserves slice-specific structural detail
while incorporating propagation-dependent context, ensuring anatomically
coherent and acoustically plausible B-mode appearance.

\vspace{-4mm}
\subsection{Contribution-Aware Gaussian Adaptation}
\label{sec:contribution-aware-gaussian-adaptation}

Conventional 3DGS density control relies on image-space gradients for
densification and opacity for pruning, which are poorly suited to ultrasound
because speckle can distort gradient cues and B-mode signals are formed by
additive echo accumulation rather than optical alpha compositing.
To overcome these limitations, we instantiate the adaptation operator
$\mathcal A$ in ULGSplat as a contribution-aware strategy consisting of
residual-guided densification and path-aware pruning, governed by the candidate
sets $\mathcal D$ and $\mathcal P$, respectively. It densifies primitives according
to their contribution to under-reconstructed pixels and removes them only when
their local and global utility cues are simultaneously weak.

\subsubsection{Residual-Guided Densification}

To identify Gaussians covering under-reconstructed regions, we define a
densification score by coupling pixel-level reconstruction errors with
Gaussian-level local contributions. Specifically, with the pixel residual
$e_{ip}
=
|\widehat I^{\mathrm{us}}_{ip}-I^{\mathrm{us}}_{ip}|$
and per-Gaussian local contribution
$c_{igp}
=
\alpha_g\eta_{ig}\nu_{ig}\phi_{igp}$,
we accumulate the local contribution
$C_g=\sum_p c_{igp}$ and its residual-weighted counterpart
$R_g=\sum_p c_{igp}e_{ip}$. The densification score is then formulated as
\begin{equation}
\omega_g
=
\frac{R_g}{C_g+\varepsilon}
\log(1+C_g),
\label{eq:contribution-aware-densification}
\end{equation}
where $\omega_g$ is the densification score of Gaussian $g$.
A larger $\omega_g$ indicates that Gaussian $g$ contributes more strongly to
under-reconstructed regions and should therefore receive additional
representation capacity.

Given the high-score quantile $Q_{q_{\mathrm d}}(\{\omega\})$, minimum score
threshold $\tau_{\omega}$, and per-event candidate budget $B$, we determine
the densification threshold and candidate set as
\begin{equation}
\tau_{\mathrm d}
=
\max
\left[
Q_{q_{\mathrm d}}(\{\omega\}),
\tau_{\omega},
\tau_{B}
\right],
\qquad
\mathcal D
=
\left\{
g
\;\middle|\;
\omega_g\geq\tau_{\mathrm d}
\right\},
\label{eq:densification-selection}
\end{equation}
where $\tau_B$ is the score threshold induced by $B$, approximately limiting
the number of selected Gaussians. The selected Gaussians are adapted according
to their physical scale: large primitives are refined by spawning child
Gaussians to subdivide broad support, whereas smaller candidates are cloned to
increase local representation capacity. The newly generated Gaussians inherit
the echo and transmission parameters of their parents and are jointly refined
during subsequent optimization.

\subsubsection{Path-Aware Pruning}

To prune Gaussians that are dispensable to both local echo formation and
global transmission, we jointly evaluate their accumulated local contribution
$C_g$ and global path contribution $A_g$, which measures the transmission
reduction induced by Gaussian $g$ over the current adaptation interval
$\mathcal K$:
\begin{equation}
A_g
=
\sum_{i\in\mathcal K}
\sum_{p:\,g\in\mathcal N^{\mathrm{glo}}_{ip}}
\left(
1-T_{igp}
\right).
\label{eq:global-path-contribution}
\end{equation}

Given the local echo-amplitude threshold $\tau_{\alpha}$, low local-contribution
quantile $Q_{q_{\mathrm p}}(\{C\})$, and low global-contribution quantile
$Q_{q_{\mathrm A}}(\{A\})$, the pruning candidate set is defined as
\begin{equation}
\mathcal P
=
\left\{
g
\;\middle|\;
\alpha_g<\tau_{\alpha},\;
C_g<Q_{q_{\mathrm p}}(\{C\}),\;
A_g<Q_{q_{\mathrm A}}(\{A\})
\right\},
\label{eq:path-aware-pruning}
\end{equation}
where $q_{\mathrm p}$ and $q_{\mathrm A}$ control the local- and
global-contribution pruning levels, respectively. 
A Gaussian is removed only when it is consistently uninformative, thereby
preserving primitives that remain relevant to ultrasound rendering. 

Through residual-guided capacity allocation and contribution-aware pruning, the
proposed adaptation regulates the Gaussian population according to
ultrasound-specific rendering demands, expanding representation capacity where
needed while suppressing redundant primitives.